Day-ahead photovoltaic (PV) forecasts for individual rooftops support household energy management, aggregator bidding and distribution-grid operation, but reported accuracy gains rarely show how close they are to a limit. We diagnose that limit for 33 residential rooftops in Belgium at 15-min resolution over 26 months, with operational numerical weather prediction (NWP) issued one day ahead. Four tests agree. Adding the measured output of every other rooftop, a proxy for perfect regional weather knowledge, raises $R^2$ by only 0.023. Residual autocorrelation falls from 0.72 at 15 min to 0.015 at 24 h, the signature of displaced cloud rather than correctable bias. Irradiance forecast error varies by 3.3\% across rooftops spread over $180 \times 210$ km. Model output statistics that remove 73\% of the systematic irradiance bias change forecast power by less than 0.0005 in $R^2$. Two methodological results generalise. Scoring the forecast against ERA5, a reanalysis from the same centre, understates its RMSE by a factor of 1.69 and reverses the sign of its bias. Hyperparameters selected on a chronological split transfer poorly to unseen rooftops, and selecting them against held-out systems halves the failure rate. An unseen rooftop receives median $R^2$ 0.704, against 0.811 pooled over rooftops seen in training, and transfer failure follows installation metadata: the observed-to-nameplate capacity ratio correlates with forecast bias at +0.709.
